\begin{textsample}{POS Dim 1 – rn – Score 28.00 – rn/rn\_em\_11.txt}  \label{ex:f1_pos_237}
3.
MODALIDADES DO ENSINO MÉDIO POTIGUAR O Ensino Médio, etapa conclusiva da Educação Básica, oportuniza aos estudantes a ampliação, aprofundamento e desenvolvimento de habilidades e competências pautadas por saberes iniciados no Ensino Fundamental.
A missão do Estado é ofertar ensino médio de qualidade e garantir a aprendizagem em todos os níveis, etapas e modalidades da educação.
Nessa etapa, os jovens, independentemente de suas origens étnico-raciais, sociais e territoriais, devem ser \textbf{vistos} e ouvidos dentro do ambiente escolar, de maneira \textbf{que} possam se (re)conhecer como parte integrante de \textbf{uma} sociedade em constante mudança.
Eles têm a oportunidade de buscarem a afirmação de suas identidades a partir do momento em \textbf{que} passam a ser protagonistas, experienciando e escrevendo a narrativa de suas histórias de vida, \textbf{uma} história em meio às reflexões sobre si mesmos e sobre os outros, marcando tomadas de posições, decisões e atitudes em prol do bem-estar e do meio social.
Para Carrano, > [... ] \textbf{uma} das mais importantes \textbf{tarefas} das instituições educativas seria de contribuir para \textbf{que} os jovens pudessem realizar escolhas conscientes sobre suas trajetórias pessoais para constituir autonomamente seus próprios acervos de valores e conhecimentos ( 2011, p. 02 ).
Nesse sentido, esse período escolar, ao momento em \textbf{que} busca contribuir com conhecimentos e estratégias para \textbf{que} os jovens continuem estudando, trará, também, elementos \textbf{que} contribuirão para as escolhas profissionais \textbf{que} os conduzam a \textbf{uma} atuação crítica, autônoma e ética na vida em sociedade.
O Estado do Rio Grande do Norte, objetivando garantir a \textbf{democratização} do Ensino Médio, em cumprimento à Lei nº 9.394/1996, oferta a educação em tempo integral e parcial, em horário diurno e noturno, em diferentes modalidades, tais como : a Educação Profissional e Tecnológica ( EPT ), a Educação Especial, a Educação Básica do Campo, a Educação Escolar Indígena, a Educação Escolar Quilombola e a Educação de Jovens e Adultos ( EJA ), adequadas às necessidades e disponibilidades das juventudes potiguares.
As ofertas mencionadas atendem \textbf{um} número de 126.859 estudantes em 329 unidades escolares, distribuídas em todo o território sob a gestão de 16 Diretorias Regionais de Ensino e Cultura ( DIRECs ), conforme representadas na figura a seguir : Figura 4 - Distribuição de escolas de Ensino Médio por DIREC Fonte : SEEC/SIGEDUC ( 2021 ).
Na relação entre o número de escolas de ensino médio da Rede Estadual e suas modalidades, observada no gráfico abaixo, prevalecem as regulares em tempo parcial, no entanto, verifica -se a expansão das de tempo integral.
Outras ofertas de representação significativa são da EJA, do noturno e as escolas \textbf{que} se destinam à EPT.
Cabe destacar o número de escolas piloto implementado ( 2019-2021 ) pelo Programa de Apoio ao Novo Ensino Médio, instituído pela Portaria MEC nº 649/2018, incluindo 100 unidades escolares do Rio Grande do Norte.
Dessas, 76 com tempo parcial ( 3.000 horas ) e 24 com tempo integral ( 4.500 horas ).
A Educação do Campo, a Educação Indígena e a Educação Especial têm, em número, unidades menos representativas diante do universo de escolas da rede.
É importante destacar \textbf{que}, embora não existam escolas exclusivas em número suficiente para assistir às comunidades do campo, indígenas e nenhuma para quilombolas, os seus estudantes são atendidos por escolas regulares, em sua maioria localizadas no meio urbano.
Gráfico 3 - Modalidades/ofertas de Ensino Médio Potiguar Fonte : SEEC/SIGEDUC ( 2021 ).
Ressaltamos \textbf{que} a somatória dos quantitativos apresentados no gráfico 03 não representam, necessariamente, o total de unidades escolares, pois pode existir mais de \textbf{uma} modalidade/oferta em \textbf{uma} mesma unidade.
As modalidades específicas têm legislação e fundamentação próprias, \textbf{que} normatizam e orientam as ofertas para a organização do trabalho pedagógico, conforme apresentaremos a seguir. 3.1 EDUCAÇÃO PROFISSIONAL E TECNOLÓGICA A Educação Profissional e Tecnológica tem sua importância no universo juvenil \textbf{enquanto} política pública nacional, pois possibilita aos jovens tanto a sua inserção no mundo do trabalho de forma \textbf{qualificada}, quanto constrói base de conhecimentos científicos, \textbf{ancorada} na ética, servindo de suporte para o desenvolvimento de projetos de intervenção social.
Nela, os \textbf{sujeitos} do Ensino Médio têm a oportunidade de atuarem com criticidade e prosseguirem em seus estudos, com protagonismo, participando de processos produtivos.
Nesse sentido, cumprindo os objetivos da educação nacional para o Ensino Médio, previstos na LDB, a Educação Profissional e Tecnológica se articula com o Ensino Médio em suas diversas modalidades, níveis e com as dimensões do trabalho, da tecnologia, da ciência e da cultura.
As Diretrizes Nacionais para a Educação Profissional Técnica de Nível Médio, consoante o art. 7º, dispõem \textbf{que} a oferta dos cursos de Educação Profissional Técnica de Nível Médio podem se apresentar de forma Articulada integrada com matrícula única na mesma instituição de ensino ou Concomitante com matrículas distintas, na mesma instituição de ensino ou em instituições diversas, mediante Convênios ou acordos de intercomplementaridade, visando ao planejamento e ao desenvolvimento de \textbf{um} projeto pedagógico unificado e integrado, e na forma Subsequente, exclusiva para quem tenha concluído o Ensino Médio.
A rede de ensino do Estado do Rio Grande do Norte atende a Educação Profissional e Tecnológica ( EPT ) com os cursos de Formação Inicial e Continuada ou Qualificação Profissional para o Trabalho e a Educação Profissional Técnica de Nível Médio, conforme o Art. nº 39, parágrafos I e II, da Lei nº 9.394/96 ( LDB ), alterada pela Lei nº 11.741/2008.
Além dos referidos cursos, as instituições \textbf{que} ofertam a Educação Profissional podem ministrar cursos especiais, estendidos à comunidade.
A construção e organização de cursos e programas de Educação Profissional, \textbf{baseados} por itinerários formativos, são dispostos por eixos tecnológicos \textbf{amparados} pelo Catálogo Nacional de Cursos Técnicos de Nível Médio, o Guia Nacional de Cursos de Formação Inicial e Continuada, mantidos pelos órgãos próprios SETEC/MEC e a Classificação Brasileira de Ocupações ( CBO ).
A EPT se assenta na integração entre educação básica e formação profissional, com a materialização de \textbf{um} currículo integrado, de forma \textbf{que} os saberes advindos da Base Comum se articulem com os da área técnica, em consonância com as dimensões do trabalho, da tecnologia, da ciência e da cultura.
Nesse sentido, a Educação Profissional não deverá ser concebida como a preparação restrita para o mercado de trabalho.
Esta, envolve, necessariamente, \textbf{uma} relação de proximidade peculiar com os caminhos históricos percorridos por nossa sociedade.
Numa perspectiva de formação integrada, os conceitos/conhecimentos devem ser apreendidos como sistema de relações de \textbf{uma} \textbf{totalidade}, em \textbf{que} o estudante terá acesso não só aos conhecimentos científicos e tecnológicos, mas também ao desenvolvimento contínuo e progressivo de habilidades e competências, de tal maneira a posicionar -se diante das questões da sociedade \textbf{contemporânea} com consciência, autonomia e criticidade.
Nos últimos anos, a sociedade \textbf{contemporânea} tem passado por mudanças significativas nos processos de produção e consumo com o avanço das novas tecnologias, como a automação, a robótica, a inteligência artificial, enfim, a realidade construída pelo universo dos algoritmos, fruto da combinação entre ciência e tecnologia \textbf{que} tem provocado inovações \textbf{que} impactam na atividade laboral e \textbf{exigem} aperfeiçoamento constante devido a rapidez em \textbf{que} elas ocorrem.
Desse modo, compartilhamos da compreensão da ‘ modernidade líquida ’, proposta por Zigmunt Bauman ( 2001 ), \textbf{que} torna evidente o estado de fluidez e de vulnerabilidade do mundo do trabalho, o \textbf{que} leva o indivíduo à necessidade de (re)agir, \textbf{face} às condições de temporalidade e transitoriedade nas relações sociais.
No plano local, observamos nas cidades potiguares fenômenos mundiais da informalidade, da terceirização e da uberização associados à realização do trabalho online, podendo ser realizado em casa, conduzindo a novas condições na economia global, regional e local e, por consequência, no mundo do trabalho.
Nessa \textbf{direção}, a preparação do estudante na Educação Profissional e Tecnológica compreende o compartilhamento de aspectos técnicos, a compreensão do indivíduo como responsável pelo seu planejamento, a capacidade de avaliação e o aperfeiçoamento do seu desempenho profissional, articulando seus saberes para responder criativamente aos desafios profissionais na sociedade \textbf{atual}.
Tais aspectos \textbf{aproximam} nosso estado de \textbf{uma} ‘ modernização ’ rápida, \textbf{que} impõe sobre os processos de aprendizagem o desenvolvimento de novas estratégias pedagógicas.
Entretanto, todo o avanço tecnológico, ainda \textbf{que} pareça meritório e útil para a facilitação de \textbf{tarefas} na vida em sociedade, é apenas meio. \textbf{Face} a essa verdade, resta \textbf{que} o desenvolvimento integral dos nossos jovens é precípuo, pois é ele \textbf{que} poderá definir o \textbf{que} fazer, como fazer e quando fazer, o \textbf{que} \textbf{influenciará} com seus atos as ações de outrem.
É essa a perspectiva da \textbf{integralidade} buscada na EPT.
Na Educação Profissional e Tecnológica a visão humanista prevalece, principalmente, porque pensa o trabalho assumido como princípio educativo e formativo.
Ou seja, o trabalho compreendido como princípio educativo é tratado como \textbf{um} aspecto \textbf{inerente} à condição humana, \textbf{uma} vez \textbf{que} o ser humano encontra -se integrado organicamente à natureza em relação de interdependência, não se \textbf{restringindo} apenas à função \textbf{didático-pedagógica} do aprender fazendo.
O Estado do Rio Grande do Norte promove a Educação Profissional \textbf{fomentada} no desenvolvimento socioeconômico-ambiental e cultural dos territórios, a partir dos seus arranjos socioprodutivos e de suas demandas.
A oferta de EPT deverá estimular e apoiar processos educativos \textbf{que} promovam a geração de trabalho e renda, assim como a emancipação do cidadão, de forma a consolidar e fortalecer os arranjos produtivos, sociais e culturais locais, a partir do mapeamento das suas potencialidades.
Para tanto, no nosso Estado, o desenvolvimento curricular prima pela flexibilização e \textbf{interdisciplinaridade}, articulando \textbf{teoria} e prática profissional de forma contextualizada e indissociável à historicidade dos envolvidos no processo de ensino, considerando os saberes científicos e os saberes desenvolvidos pelos \textbf{sujeitos} em outros espaços sociais. \textbf{Logo}, a sua proposta pedagógica para a Educação Profissional reconhece a integração sistêmica entre a exploração sustentável dos recursos naturais, o desenvolvimento tecnológico e a mudança social, \textbf{que} \textbf{implicam} em atitudes éticas, construções de propostas tecnológicas e processos de ensino e aprendizagem mais significativos.

% matched lemmas: amparar, ancorar, aproximar, atual, basear, contemporâneo, democratização, didático-pedagógico, direção, enquanto, exigir, face, fomentar, implicar, inerente, influenciar, integralidade, interdisciplinaridade, logo, qualificar, que, restringir, sujeito, tarefa, teoria, totalidade, um, ver
\end{textsample}
